\smallskip\noindent\textbf{Sliders along real-music axes.} Seven axes were turned into sliders with the set trainer, the sets being the top and bottom 20\% of 24 seeds per prompt along the axis (Table~\ref{tab:axes}). All \axN{} follow their axis, with rank correlations between \axRhoLo{} and \axRhoHi{}. The arousal slider moves one prompt and seed 1.19 standard deviations of real music along its axis, more than the standard deviation of the model's whole output over 648 prompts (0.83) and 2.1 times the standard deviation that seeds of one prompt produce. Similarity to the unsteered clip is \axKeepLo{}--\axKeepHi{} at $\pm1$, enjoyment is at or above the unsteered level, and musicality drops by at most 0.14.

\begin{table}[t]\centering\footnotesize\setlength{\tabcolsep}{2.5pt}
\begin{tabular}{llrrrrrrr}
\toprule
Axis & From & $\rho$ & Ord. & Real & Seeds & Kept & CE & Mus. \\
\midrule
acoustic/electr. & sets & 0.55 & 0.93 & 0.64 & 1.48 & 0.86 & 7.11 & 2.68 \\
arousal & sets & 0.82 & 0.99 & 1.19 & 2.14 & 0.84 & 7.13 & 2.70 \\
classical/funk & sets & 0.64 & 0.96 & 0.77 & 1.65 & 0.85 & 7.12 & 2.69 \\
jazz/electronic & sets & 0.61 & 0.94 & 0.61 & 1.23 & 0.87 & 7.09 & 2.70 \\
piano & sets & 0.63 & 0.96 & 0.75 & 1.42 & 0.87 & 7.08 & 2.69 \\
strings/synth & sets & 0.40 & 0.82 & 0.44 & 0.81 & 0.87 & 7.09 & 2.69 \\
valence & sets & 0.50 & 0.93 & 0.62 & 1.23 & 0.87 & 7.07 & 2.73 \\
\bottomrule
\end{tabular}
\caption{Sliders along axes of real music on ACE-Step, positions $-1$ to $+1$. From: a prompt pair made of the axis's tags, or two sets of the model's clips. $\rho$, Ord.: rank correlation between position and the output's projection on the axis, and share of trajectories with $-1$ and $+1$ in the right order. Real, Seeds: movement along the axis between $-1$ and $+1$ in standard deviations of real recordings and of one prompt's seeds. Kept: CLAP similarity to the unsteered clip. CE, Mus.: enjoyment and SongEval musicality at $\pm1$ (unsteered: 6.95 and 2.77).}\label{tab:axes}
\end{table}


\smallskip\noindent\textbf{Each trainer moves what it was trained on.} Sorting clips by a waveform descriptor, residualised against loudness, brightness, and enjoyment, gives sliders that are selective. Over positions $-1$ to $+1$ the set-trained harmony slider reaches $\rho=0.72$ on harmonic change rate with selectivity 6.6, against 0.27 and 0.8 for the prompt-pair slider; density and ensemble show 4.7 and 2.4 against 1.5 and 1.5. The embedding tells the opposite story: the MuQ-MuLan text direction for harmony follows the prompt-pair slider ($\rho=0.63$) and ignores the set-trained one ($-0.07$). Neither kind of score can stand in for the other. Set-trained sliders stop at $\pm1$: beyond it similarity falls below 0.72 and musicality drops by 0.3--0.5 in both directions while CE stays flat, a failure only the second predictor shows. Tempo did not become a slider this way ($\rho=0.06$).
